\documentclass[times]{san}

\bibliografia{bibliografia.bib}

\uczelnia{Społeczna Akademia Nauk}
\przedmiot{Systemy Szkieletowe}
\autor{Jan Kowalski}
\nralbumu{123456}
\grupa{3}
\tytul{Nazwa projektu}
\prowadzacy{mgr inż. Krystian Gumiński}
\miasto{Łódź}
\rok{2026}

\begin{document}

\maketitlepage
\spistresci

\section{Wstęp}

Celem niniejszego rozdziału jest ogólne wprowadzenie do tematyki projektu. Praca dotyczy zagadnień związanych z projektowaniem i implementacją aplikacji informatycznej. W dalszej części omówiono cel projektu, zastosowaną architekturę oraz kluczowe aspekty implementacyjne \cite{knuth1984}.

\section{Cel projektu}

Celem projektu jest określenie problemu, który rozwiązuje projektowana aplikacja, oraz zdefiniowanie celów, jakie ma ona spełniać. Obejmują one zarówno wymagania funkcjonalne, jak i niefunkcjonalne, takie jak wydajność, skalowalność oraz użyteczność.

\section{Przegląd istniejących rozwiązań}

W niniejszym rozdziale przedstawiono analizę trzech podobnych rozwiązań dostępnych na rynku. Na rysunku \ref{fig:lorem} zaprezentowano przykładowy interfejs jednego z analizowanych systemów, który posłużył jako punkt odniesienia w dalszych pracach projektowych \cite{website_example}.

\begin{figure}[h]
\centering
\includegraphics[width=0.6\textwidth]{psyduck.jpg}
\caption{Przykładowy interfejs analizowanego systemu}
\label{fig:lorem}
\end{figure}

\subsection{Rozwiązanie A}
Lorem ipsum dolor sit amet, consectetur adipiscing elit. Curabitur eget felis justo. Duis ut imperdiet lorem, condimentum efficitur nunc. Praesent lobortis pharetra gravida. Ut tempus, metus a mattis ultrices, leo urna placerat turpis, quis lobortis enim diam at mi. Aenean rutrum ex eu ultricies consequat. In ornare nisl quis urna convallis suscipit. Ut venenatis, neque vitae hendrerit gravida, nisi mauris pharetra ex, sit amet varius felis mi hendrerit risus. Fusce luctus, augue vel pretium faucibus, ex nulla tempus ante, bibendum laoreet arcu ex et eros. Aenean cursus neque iaculis ligula euismod, nec laoreet velit tristique. Ut posuere, ante id gravida posuere, ante ex sodales ipsum, at vulputate elit metus at nunc.

\subsection{Rozwiązanie B}
Vivamus sagittis lacus sapien, non suscipit velit mollis at. Ut quis eros condimentum, imperdiet lacus sit amet, blandit est. Nullam dapibus ante id sapien mollis dignissim. Vestibulum magna leo, pulvinar non nisl non, posuere hendrerit libero. In id consectetur tortor, pretium blandit sem. Mauris in justo eu justo consectetur euismod. Ut eu auctor tortor. Ut nec orci in massa lobortis commodo eget et nisi. Pellentesque aliquam lectus tellus, et consequat magna aliquam eget. Quisque tincidunt mattis libero id faucibus. Integer gravida ex finibus, venenatis quam ut, cursus justo. Suspendisse vestibulum, eros non hendrerit semper, ante mauris porttitor est, vel consectetur enim ipsum et lorem.

\subsection{Rozwiązanie C}
Donec efficitur purus non convallis tempus. Etiam orci erat, rhoncus nec massa nec, efficitur ullamcorper nunc. Sed lobortis, erat vitae pulvinar accumsan, turpis arcu condimentum velit, in suscipit purus neque id elit. 

Nullam varius est lorem, eu dictum justo faucibus ac. Phasellus purus metus, consequat quis bibendum eget, sodales a nibh. Duis ornare justo velit, non sagittis tortor egestas ac. Nulla facilisi. Mauris vitae imperdiet quam. Maecenas vitae turpis sed est iaculis volutpat. 

Pellentesque fermentum dui nec dolor gravida ullamcorper. Praesent fermentum auctor elit, ac scelerisque sapien dictum a. Morbi posuere libero sit amet quam feugiat condimentum. Cras rutrum mi ligula, eget elementum urna egestas in. Maecenas vel arcu eget lorem viverra luctus.



\section{Opis architektury}

Rozdział zawiera szczegółowy opis architektury projektowanego systemu,
z uwzględnieniem podziału na warstwy oraz sposobu komunikacji pomiędzy
poszczególnymi komponentami. Tabela \ref{tab:lorem} przedstawia zestawienie
głównych elementów architektury wraz z ich rolą w systemie.

\begin{table}[h]
\centering
\caption{Główne komponenty architektury systemu}
\label{tab:lorem}
\begin{tabular}{|c|l|c|}
\hline
Lp. & Komponent          & Warstwa \\
\hline
1   & Backend API        & Logika   \\
2   & Frontend Web       & Prezentacja \\
3   & Baza danych        & Dane     \\
\hline
\end{tabular}
\end{table}

\section{Opis implementacji}

W rozdziale omówiono kluczowe aspekty implementacyjne projektu.
Opis został podzielony na trzy części obejmujące backend, frontend
oraz bazę danych. Przedstawiono zastosowane technologie, strukturę
kodu oraz istotne decyzje projektowe.

\subsection{Frontend}
Donec efficitur purus non convallis tempus. Etiam orci erat, rhoncus nec massa nec, efficitur ullamcorper nunc. Sed lobortis, erat vitae pulvinar accumsan, turpis arcu condimentum velit, in suscipit purus neque id elit. Nullam varius est lorem, eu dictum justo faucibus ac. Phasellus purus metus, consequat quis bibendum eget, sodales a nibh. Duis ornare justo velit, non sagittis tortor egestas ac. Nulla facilisi. Mauris vitae imperdiet quam. 

Maecenas vitae turpis sed est iaculis volutpat. Pellentesque fermentum dui nec dolor gravida ullamcorper. Praesent fermentum auctor elit, ac scelerisque sapien dictum a. Morbi posuere libero sit amet quam feugiat condimentum. Cras rutrum mi ligula, eget elementum urna egestas in. Maecenas vel arcu eget lorem viverra luctus.

\subsection{Backend}
Donec efficitur purus non convallis tempus. Etiam orci erat, rhoncus nec massa nec, efficitur ullamcorper nunc. Sed lobortis, erat vitae pulvinar accumsan, turpis arcu condimentum velit, in suscipit purus neque id elit. Nullam varius est lorem, eu dictum justo faucibus ac. Phasellus purus metus, consequat quis bibendum eget, sodales a nibh. Duis ornare justo velit, non sagittis tortor egestas ac. Nulla facilisi. Mauris vitae imperdiet quam. 

Maecenas vitae turpis sed est iaculis volutpat. Pellentesque fermentum dui nec dolor gravida ullamcorper. Praesent fermentum auctor elit, ac scelerisque sapien dictum a. Morbi posuere libero sit amet quam feugiat condimentum. Cras rutrum mi ligula, eget elementum urna egestas in. Maecenas vel arcu eget lorem viverra luctus.

\subsection{Baza Danych}
Donec efficitur purus non convallis tempus. Etiam orci erat, rhoncus nec massa nec, efficitur ullamcorper nunc. Sed lobortis, erat vitae pulvinar accumsan, turpis arcu condimentum velit, in suscipit purus neque id elit. Nullam varius est lorem, eu dictum justo faucibus ac. Phasellus purus metus, consequat quis bibendum eget, sodales a nibh. Duis ornare justo velit, non sagittis tortor egestas ac. Nulla facilisi. Mauris vitae imperdiet quam. 

Maecenas vitae turpis sed est iaculis volutpat. Pellentesque fermentum dui nec dolor gravida ullamcorper. Praesent fermentum auctor elit, ac scelerisque sapien dictum a. Morbi posuere libero sit amet quam feugiat condimentum. Cras rutrum mi ligula, eget elementum urna egestas in. Maecenas vel arcu eget lorem viverra luctus.

\subsection{Docker}
Donec efficitur purus non convallis tempus. Etiam orci erat, rhoncus nec massa nec, efficitur ullamcorper nunc. Sed lobortis, erat vitae pulvinar accumsan, turpis arcu condimentum velit, in suscipit purus neque id elit. Nullam varius est lorem, eu dictum justo faucibus ac. Phasellus purus metus, consequat quis bibendum eget, sodales a nibh. Duis ornare justo velit, non sagittis tortor egestas ac. Nulla facilisi. Mauris vitae imperdiet quam. 

Maecenas vitae turpis sed est iaculis volutpat. Pellentesque fermentum dui nec dolor gravida ullamcorper. Praesent fermentum auctor elit, ac scelerisque sapien dictum a. Morbi posuere libero sit amet quam feugiat condimentum. Cras rutrum mi ligula, eget elementum urna egestas in. Maecenas vel arcu eget lorem viverra luctus.

\subsection{Testy}
Donec efficitur purus non convallis tempus. Etiam orci erat, rhoncus nec massa nec, efficitur ullamcorper nunc. Sed lobortis, erat vitae pulvinar accumsan, turpis arcu condimentum velit, in suscipit purus neque id elit. Nullam varius est lorem, eu dictum justo faucibus ac. Phasellus purus metus, consequat quis bibendum eget, sodales a nibh. Duis ornare justo velit, non sagittis tortor egestas ac. Nulla facilisi. Mauris vitae imperdiet quam. 

Maecenas vitae turpis sed est iaculis volutpat. Pellentesque fermentum dui nec dolor gravida ullamcorper. Praesent fermentum auctor elit, ac scelerisque sapien dictum a. Morbi posuere libero sit amet quam feugiat condimentum. Cras rutrum mi ligula, eget elementum urna egestas in. Maecenas vel arcu eget lorem viverra luctus.

\subsection{Monitoring}
Donec efficitur purus non convallis tempus. Etiam orci erat, rhoncus nec massa nec, efficitur ullamcorper nunc. Sed lobortis, erat vitae pulvinar accumsan, turpis arcu condimentum velit, in suscipit purus neque id elit. Nullam varius est lorem, eu dictum justo faucibus ac. Phasellus purus metus, consequat quis bibendum eget, sodales a nibh. Duis ornare justo velit, non sagittis tortor egestas ac. Nulla facilisi. Mauris vitae imperdiet quam. 

Maecenas vitae turpis sed est iaculis volutpat. Pellentesque fermentum dui nec dolor gravida ullamcorper. Praesent fermentum auctor elit, ac scelerisque sapien dictum a. Morbi posuere libero sit amet quam feugiat condimentum. Cras rutrum mi ligula, eget elementum urna egestas in. Maecenas vel arcu eget lorem viverra luctus.

\section{Podsumowanie i wnioski}

W podsumowaniu przedstawiono rezultaty zrealizowanej pracy,
napotkane problemy oraz sposoby ich rozwiązania. Wskazano również
możliwe kierunki dalszego rozwoju aplikacji.

\spistabel
\spisrysunkow
\printbibliography

\end{document}